\documentclass[11pt]{article}
\usepackage[margin=1in]{geometry}
\usepackage{times}
\usepackage[hidelinks]{hyperref}
\usepackage{enumitem}
\setlist[itemize]{topsep=4pt, itemsep=2pt, parsep=0pt}

\title{\textbf{AI Overlay Assistant for Coding Help and Meeting Notes}}
\author{}
\date{}

\begin{document}
\pagestyle{plain}
\thispagestyle{empty}

\begin{center}
{\LARGE \textbf{AI Overlay Assistant for Coding Help and Meeting Notes}}\\[6pt]
Project Proposal for UCS503P\\
Submitted to: Ms. Paramveer Kaur\\
Thapar Institute of Engineering and Technology\\[10pt]
Sangam Sangwan, CSED \quad Samyak Jain, CSED
\end{center}

\vspace{10pt}

\section{Higher-Order Goal}
This project aims to save students time in two ordinary situations: getting stuck on a
coding problem, and writing up notes after a meeting. The broader motivation is to reduce
small, recurring time losses in day-to-day academic and project work. The immediate
engineering objective is a desktop overlay assistant that (a) gives a short hint when a
student is stuck on a coding problem on their own screen, and (b) turns a meeting's audio
into a short text summary.

\section{Time-to-Value}
\begin{itemize}
  \item \textbf{We will build one feature at a time:} the coding-hint feature will be built
    and tested first, since it needs a simpler capture path (a single screenshot) than
    continuous meeting audio.
  \item \textbf{We will start with the simplest working capture:} a hotkey that grabs the
    current screen, before attempting continuous or background capture.
  \item \textbf{We will keep the first version small:} a single working flow (trigger a
    hint, get a hint back) is the target for the first iteration, not a finished product.
\end{itemize}

\section{Problem Statement}
Two small but common problems motivate this project:
\begin{itemize}
  \item Coding: when a student gets an error or gets stuck, they usually search the web or
    a forum for help. This works, but takes time and often returns answers that don't match
    their exact code.
  \item Meetings: after a project meeting, someone has to write up what was discussed and
    what needs to be done. This is often skipped or done incompletely, since it takes extra
    time after the meeting is already over.
\end{itemize}
Neither problem is severe, but both are frequent enough that a small tool addressing them is
a reasonable, testable engineering project.

\section{Proposed Solution}

\subsection{Overview}
A desktop overlay assistant with two independent modes:
\begin{itemize}
  \item Coding help: while working on a problem, the student triggers the overlay with a
    hotkey; it captures the current screen and returns a short hint (not a full solution).
  \item Meeting notes: the student turns on audio capture for a meeting; the assistant
    transcribes it and produces a short summary with any action items once the meeting ends.
\end{itemize}
Both modes share our own backend for hint and summary generation and logging. The overlay
client itself --- the floating panel, screen capture, and audio capture --- is built by us
for this project.

\subsection{Core Workflow}
Coding help:
\begin{itemize}
  \item Student presses a hotkey while stuck; the overlay captures the current screen.
  \item Our backend sends the captured content to an LLM API with a prompt that asks for a
    short hint, not the full answer.
  \item Student marks whether the hint was useful.
\end{itemize}
Meeting notes:
\begin{itemize}
  \item Student turns on audio capture (microphone and system audio) for a meeting.
  \item Audio is transcribed while capture is on; our backend builds a running summary.
  \item At the end of the session, the student receives a short summary with any action
    items, which they can view and download.
\end{itemize}

\subsection{Operational Constraints}
\begin{itemize}
  \item For this project, both continuous capture and any hidden-display behaviour are
    scoped to personal, individual use --- a student's own coding practice and their own
    meetings --- not proctored exams, interviews, or another party's call without their
    knowledge.
  \item The overlay keeps a visible on/off indicator and an explicit consent step before any
    capture starts.
  \item Usable on a normal laptop; requires OS-level screen-recording and microphone
    permissions.
  \item Built using existing, off-the-shelf APIs for language and speech processing with a
    bring-your-own API key; the overlay window and capture logic are our own implementation.
\end{itemize}

\section{Solution Approach}
This is an engineering course project, so the focus is on building a working, well-tested
tool using existing components for language and speech processing, not on research or new
algorithms.

\subsection{Hint Generation Approach}
\begin{itemize}
  \item Hints are generated by calling an existing large language model API with a prompt
    that asks for a short hint, not a complete solution.
  \item This is a deliberate, simple design choice --- it keeps the tool useful for learning
    rather than for skipping the problem --- and is implemented with prompt wording and
    length limits, not new algorithms.
\end{itemize}

\subsection{System Architecture}
A client-server architecture, built for this project:
\begin{itemize}
  \item Client: our own Electron-based overlay (floating panel, screen capture, audio
    capture, and an optional hidden-display mode).
  \item Backend API: our own service that receives captured screen content or audio, calls
    the LLM and speech-to-text APIs, and returns hints or summaries.
  \item Database: stores users, hint/summary logs, and consent records for evaluation.
\end{itemize}

\subsection{CI/CD Approach}
CI/CD will be used to:
\begin{itemize}
  \item Automatically build and run tests on every change, for both the client and the
    backend.
  \item Produce repeatable deployment artifacts to staging.
  \item Enable safe and frequent releases of incremental features.
\end{itemize}

\section{Evaluation Criteria}

\subsection{Primary Evaluation Metric}
Time-to-hint (TTH): time between a student triggering a hint request and receiving a hint in
the overlay.
\begin{itemize}
  \item Target: median TTH $\leq$ 1 minute during pilot.
  \item Attribution: measured from client/server timestamps (trigger vs. hint returned).
\end{itemize}

\subsection{Secondary Evaluation Metrics}
\begin{itemize}
  \item Hint usefulness: percentage of hints marked useful by students (simple thumbs
    up/down).
  \item Summary usefulness: student rating (1--5) of whether the meeting summary captured
    the discussion accurately.
  \item Reliability: uptime of the staging/production backend during business hours during
    the pilot.
\end{itemize}

\subsection{Pilot Validation Plan}
\begin{itemize}
  \item Ask 8--10 classmates to use the coding-help mode on their own real assignment
    problems over one to two weeks.
  \item Ask 2--3 project groups to try the meeting-notes mode on their own meeting.
  \item Collect the ratings above and a few short comments to guide the next iteration.
\end{itemize}

\section{Scalability}
1. \textbf{Scalable (theoretically):} the backend is stateless and decoupled from the
client, so more users can be served by running more backend instances; database queries for
past hints/summaries are indexed and paginated.

2. \textbf{Operationally deployable (practically):} the backend runs on standard web hosting
with a managed database, a containerized service, and a CI/CD pipeline for staging and
production.

\section{Technology Stack}
A mature set of tooling is available to implement this project:
\begin{itemize}
  \item A desktop application framework (e.g. Electron) for building the overlay client.
  \item Web framework for our own backend REST API.
  \item Managed relational database.
  \item An existing LLM API for hint and summary generation.
  \item An existing speech-to-text API for transcribing meeting audio.
  \item Authentication approach (token-based or session-based).
  \item Test framework for unit/integration tests.
  \item CI runner and staging deployment target.
\end{itemize}
Using existing APIs for language and speech processing lets us focus our own engineering
effort on the overlay client, backend, prompt design, and evaluation, rather than on model
training or infrastructure that already exists as a mature service.

\section{Project Scope and Deliverables}

\subsection{Initial Deliverable}
\begin{itemize}
  \item A minimal overlay client: floating panel plus a hotkey-triggered screenshot capture.
  \item Coding-help mode: hotkey-triggered hint request, short hint returned, marked useful
    or not.
  \item Basic user accounts, a consent step before capture starts, and session logging.
  \item Basic logging and timestamps for the TTH metric.
  \item CI: build + backend unit tests + linting.
  \item CD to staging with a single command or pipeline job.
\end{itemize}

\subsection{Subsequent Deliverables}
\begin{itemize}
  \item Continuous audio capture and the meeting-notes mode (transcript through to summary).
  \item Optional hidden-display mode for the overlay, if time allows.
  \item Simple dashboard showing hint usefulness and summary ratings over time.
  \item Minor refinements based on pilot feedback (e.g. hint length, prompt wording).
\end{itemize}

\section{Risks and Mitigations}
\begin{itemize}
  \item \textbf{Misuse in exams or interviews:} continuous capture and any hidden-display
    mode could, in principle, be used to hide AI assistance during a proctored exam or
    interview. We scope this project's use to personal practice and the student's own
    meetings only, require an explicit consent step before capture starts, and log every
    session for the user's own review.
  \item \textbf{Capture and overlay behaviour is nontrivial to build:} screen capture, audio
    capture, and hidden-window behaviour can behave differently across OS versions and are
    more work than they first appear. We scope the first iteration to a single
    hotkey-triggered screenshot, and treat continuous capture and hidden-display mode as
    later, optional deliverables rather than promising them upfront.
  \item \textbf{Over-reliance on hints:} keep hints short and incomplete by design, so the
    tool supports rather than replaces problem-solving.
  \item \textbf{Data privacy:} only process content captured during an active, consented
    session; do not store screenshots or audio longer than needed to generate a hint or
    summary; keep summaries private to the user.
  \item \textbf{API cost or downtime:} add basic rate limits per user and a fallback message
    if the LLM or speech API is unavailable.
  \item \textbf{Low pilot participation:} recruit pilot users from within our own
    course/project group first, where participation is easier to arrange.
  \item \textbf{Evaluation ambiguity:} define ``useful hint'' and ``accurate summary'' as
    simple rating questions before development begins, so results are easy to compare.
\end{itemize}

\section{Summary}
This project proposes a small, two-mode desktop overlay assistant --- a coding-hint helper
and a meeting-notes generator --- built as an original implementation using existing
language and speech APIs rather than new research. It follows the course criteria by
shipping a narrow first iteration quickly, using CI/CD for safe releases, and evaluating
itself with simple, measurable metrics (time-to-hint and usefulness ratings). Capture and
display behaviour that is harder to build (continuous audio capture, hidden-display mode) is
deliberately phased into later deliverables, and their use is scoped to personal practice
with consent and logging built in.

\end{document}
